% !TeX root = ../深度学习与大语言模型技术原理.tex
% 由主文档通过 \input 载入；请编译主文档。
\setcounter{chapter}{3} % 保留原书章号。
\chapter{分词、自回归建模与预训练}
\label{ch:pretraining}

\section{分词器决定模型的离散世界}

模型处理的是有限词表中的编号，而非直接处理人类理解的单词。\emph{词元（token）}可以是字、子词、字节片段或特殊控制符；分词器（tokenizer）把文本映成词元序列，并提供反向解码。词表较大通常使序列较短，却增加嵌入层和输出层开销；词表较小减少每个词元的分类空间，却可能增加注意力和生成步数。不同语言、代码和数字的切分密度不同，因此词表设计影响同一上下文窗口能够容纳多少有效信息。

\emph{字节对编码（byte-pair encoding, BPE）}从较小单元开始，反复统计相邻符号对并合并选定的高频对，形成有顺序的合并规则。训练结束后，编码阶段使用这些规则切分新文本。字节级 BPE 以字节为底层单位，有助于覆盖任意可编码文本，但一个词元不一定对应完整的 Unicode 字符，也不保证数字按数学位数切分。

Unigram 分词模型为候选子词 $s$ 赋概率 $p(s)$。对一种合法切分 $z=(s_1,\ldots,s_k)$，设 $p(z)=\prod_jp(s_j)$；文本概率对全部合法切分求和，训练据此调整子词概率并裁减词表。给定词表后，最优切分可用动态规划求得。它与 BPE 的贪心合并机制不同。SentencePiece 是支持相关分词方法的软件工具，并非与 BPE、Unigram 同层次的第三种单一算法\cite{sentencepiece}。

分词器还必须规定规范化、空格、未知字符和特殊词元的处理。\emph{对话模板（chat template）}进一步把角色、消息边界、工具返回和回答起点串成模型输入。训练时使用一种模板、部署时换另一种模板，会改变条件分布；重复添加开始标记、把结束标记截断或混淆用户与助手角色，都可能导致行为异常。

\section{下一词元预测与标签移位}

对词元序列 $x_{1:n}$，链式法则给出
\begin{equation}
p_\theta(x_{1:n})=\prod_{t=1}^np_\theta(x_t\mid x_{<t}),\qquad
\mathcal L_{\mathrm{LM}}=-\sum_{t=1}^n\log p_\theta(x_t\mid x_{<t}).
\label{eq:autoregressive}
\end{equation}
这是一种概率分解，条件分布通过神经网络参数化。序列的开始由约定的起始上下文表示，结束词元也属于预测目标。若限制最大生成长度，相当于额外规定截断规则，不能把截断自动理解为模型自然结束。

训练采用\emph{教师强制（teacher forcing）}：每个位置都读取真实前缀，并行预测后一个词元。给定输入 $(\mathrm{BOS},x_1,x_2)$，对应目标是 $(x_1,x_2,\mathrm{EOS})$，而不是输入自身。因果掩码防止位置看到未来标签。推理时前缀含模型自己的输出，偏离训练前缀后会产生分布差异；这种差异解释了为何训练损失下降不自动保证长序列生成稳定。

\emph{困惑度（perplexity）}定义为 $\exp(\mathcal L_{\mathrm{LM}}/N_{\mathrm{valid}})$，其中 $N_{\mathrm{valid}}$ 为参与评分的词元总数。它是平均负对数似然的单调变换，较低说明模型对指定测试文本赋予更高概率。不同分词器改变了词元单位，因此词元级困惑度不能直接跨词表比较；不同评测掩码也会改变数值。

\begin{example}
若一段测试文本的三个目标词元条件概率分别为 $1/2,1/4,1/8$，求该段平均负对数似然与困惑度。
\end{example}
\begin{solve}
序列概率为 $1/64$，总负对数似然为 $\log64=6\log2$，平均值为 $2\log2$，所以困惑度为 $4$。该数值表示模型在这个计分单位上的平均不确定性，并不表示每一步恰好在四个词元中均匀选择。
\end{solve}

\section{嵌入、输出投影与训练目标的边界}

嵌入矩阵 $E\in\mathbb R^{V\times d}$ 把词元编号映为向量。经过 $L$ 层变换后，位置 $t$ 的隐藏向量 $h_t$ 通过输出矩阵 $W_{\mathrm{out}}\in\mathbb R^{V\times d}$ 产生 logits。权重绑定（weight tying）令 $W_{\mathrm{out}}=E$，减少参数并约束输入与输出表征；是否绑定属于架构选择，并非语言模型的必要条件。

预训练（pre-training）在广泛语料上学习条件分布，继续预训练（continued pre-training）在已有模型基础上继续语言建模，监督微调则着重学习指令条件下的目标回答。三者可以共享交叉熵形式，但数据、掩码、分布与预期行为不同。把某领域文档喂给继续预训练，未必同时教会模型按该领域问答格式回答；高质量任务示例仍可能必要。

对于稠密 Transformer，在忽略部分注意力和输出层成本、以参数矩阵乘法为主的估算下，训练浮点运算量常粗略写为 $C\approx6PN_{\mathrm{tok}}$。这来自每个参数在每个词元上的前向及反向矩阵乘法成本估计，不适用于所有序列长度、稀疏结构和重计算配置。扩大模型与增加数据如何分配预算，需要缩放实验（scaling experiment）验证；幂律拟合是指定区间内的经验模型，不能无条件外推。

